% LaTeX syntax highlighting sample
\documentclass[12pt,a4paper]{article}

% Packages
\usepackage[utf8]{inputenc}
\usepackage[T1]{fontenc}
\usepackage{amsmath,amssymb,amsthm}
\usepackage{graphicx}
\usepackage[colorlinks=true,linkcolor=blue]{hyperref}
\usepackage{listings}
\usepackage{booktabs}
\usepackage{geometry}
\usepackage{fancyhdr}
\usepackage[english]{babel}
\usepackage{xcolor}

% Page layout
\geometry{margin=1in}
\pagestyle{fancy}
\fancyhf{}
\fancyhead[L]{\leftmark}
\fancyhead[R]{\thepage}

% Custom commands
\newcommand{\R}{\mathbb{R}}
\newcommand{\N}{\mathbb{N}}
\newcommand{\norm}[1]{\left\| #1 \right\|}
\renewcommand{\vec}[1]{\boldsymbol{#1}}

% Theorem environments
\newtheorem{theorem}{Theorem}[section]
\newtheorem{lemma}[theorem]{Lemma}
\newtheorem{corollary}[theorem]{Corollary}

% Document metadata
\title{A Comprehensive \LaTeX{} Example}
\author{Jane Doe \and John Smith}
\date{\today}

\begin{document}

\maketitle
\tableofcontents

\section{Introduction}

This document demonstrates various \LaTeX{} features for syntax
highlighting purposes. We use both inline math $x^2 + y^2 = r^2$
and display math environments.

\subsection{Basic Formatting}

Text can be \textbf{bold}, \textit{italic}, \texttt{monospaced},
or \emph{emphasized}. Special characters: \& \% \$ \# \_ \{ \}

\subsection*{Mathematical Expressions}

The quadratic formula:
$$x = \frac{-b \pm \sqrt{b^2 - 4ac}}{2a}$$

An aligned equation system:
\begin{align}
    \nabla \cdot \vec{E} &= \frac{\rho}{\epsilon_0} \label{eq:gauss} \\
    \nabla \times \vec{B} &= \mu_0 \vec{J} + \mu_0 \epsilon_0 \frac{\partial \vec{E}}{\partial t}
\end{align}

Inline math: Let $f: \R \to \R$ be a continuous function. Then for
all $\epsilon > 0$, there exists $\delta > 0$ such that
$|x - x_0| < \delta$ implies $|f(x) - f(x_0)| < \epsilon$.

\section{Theorems and Proofs}

\begin{theorem}[Fundamental Theorem of Calculus]
    Let $f$ be continuous on $[a, b]$. Then
    \[
        \int_a^b f(x) \, dx = F(b) - F(a)
    \]
    where $F$ is any antiderivative of $f$.
\end{theorem}

\begin{proof}
    The proof follows from the definition of the Riemann integral
    and the Mean Value Theorem. See~\cite{rudin1976} for details.
\end{proof}

\begin{lemma}
    For all $n \in \N$, we have $\sum_{k=1}^{n} k = \frac{n(n+1)}{2}$.
\end{lemma}

\section{Tables and Figures}

\begin{table}[htbp]
    \centering
    \caption{Comparison of algorithms}
    \label{tab:algorithms}
    \begin{tabular}{lccc}
        \toprule
        Algorithm & Time & Space & Stable \\
        \midrule
        Quicksort & $O(n \log n)$ & $O(\log n)$ & No \\
        Mergesort & $O(n \log n)$ & $O(n)$ & Yes \\
        Heapsort  & $O(n \log n)$ & $O(1)$ & No \\
        \bottomrule
    \end{tabular}
\end{table}

See Table~\ref{tab:algorithms} and Figure~\ref{fig:example}.

\begin{figure}[htbp]
    \centering
    \includegraphics[width=0.8\textwidth]{figure.pdf}
    \caption{An example figure.}
    \label{fig:example}
\end{figure}

\section{Lists}

\begin{enumerate}
    \item First item with \textbf{bold text}
    \item Second item with math: $e^{i\pi} + 1 = 0$
    \item Third item
    \begin{itemize}
        \item Nested bullet
        \item Another bullet
    \end{itemize}
\end{enumerate}

\section{Code Listing}

\begin{lstlisting}[language=Python, caption=Hello World]
def hello(name):
    print(f"Hello, {name}!")

if __name__ == "__main__":
    hello("World")
\end{lstlisting}

\begin{verbatim}
This text is displayed verbatim.
No \LaTeX commands are processed here.
\end{verbatim}

\section{Conclusion}

For more information, see the documentation\footnote{Available at
\url{https://www.latex-project.org}}.

\bibliography{references}
\bibliographystyle{plain}

\end{document}
